\chapter{Time-Independent Schr\"odinger Equation}

\section*{Introduction}

The hydrogen atom has only certain allowed energy levels---not a continuous range. This is why hydrogen gas emits only specific colors of light (spectral lines). The Schr\"odinger equation predicts these energies exactly, matching experiment to ten decimal places---the most precise prediction in all of science. Every quantum device, from transistors to lasers, relies on solutions to this equation.



\[
\hat{H}\psi = E\psi \quad \text{where } \hat{H} = -\frac{\hbar^2}{2m}\nabla^2 + V(\mathbf{r})
\]

\section*{Fundamental Equations Used}

The derivation starts from the time-dependent Schr\"odinger equation using separation of variables.

\section*{Assumptions}

\textbf{Assumptions:} The potential $V(\mathbf{r})$ is time-independent. The wave function $\psi$ provides a complete description. The Hamiltonian is Hermitian (guaranteeing real energy eigenvalues).


\textbf{Limitations:} Non-relativistic (the Dirac equation is needed for relativistic speeds). Does not include spin unless the Pauli term is added. The many-body problem is intractable analytically for more than a few particles.


\section*{Mathematical Derivation}

\textbf{Step 1: The Time-Dependent Schr\"odinger Equation.}

We begin with the full time-dependent Schr\"odinger equation for a particle of mass $m$ in a potential $V(\mathbf{r})$:

\begin{equation}
i\hbar\,\frac{\partial \Psi(\mathbf{r}, t)}{\partial t} = \hat{H}\,\Psi(\mathbf{r}, t) = \left[-\frac{\hbar^2}{2m}\nabla^2 + V(\mathbf{r})\right]\Psi(\mathbf{r}, t).
\end{equation}

This is the fundamental postulate of non-relativistic quantum mechanics: it governs the evolution of the wave function $\Psi$ in time.

\textbf{Step 2: Separation of Variables.}

We seek solutions in which the spatial and temporal parts factorize. Assume:

\begin{equation}
\Psi(\mathbf{r}, t) = \psi(\mathbf{r})\,\phi(t).
\end{equation}

Substitute into the time-dependent Schr\"odinger equation:

\begin{equation}
i\hbar\,\psi(\mathbf{r})\,\frac{d\phi}{dt} = \left[-\frac{\hbar^2}{2m}\nabla^2\psi(\mathbf{r}) + V(\mathbf{r})\,\psi(\mathbf{r})\right]\phi(t).
\end{equation}

\textbf{Step 3: Separating the Variables.}

Divide both sides by $\psi(\mathbf{r})\,\phi(t)$:

\begin{equation}
i\hbar\,\frac{1}{\phi}\,\frac{d\phi}{dt} = \frac{1}{\psi}\left[-\frac{\hbar^2}{2m}\nabla^2\psi + V\,\psi\right].
\end{equation}

The left side depends only on $t$; the right side depends only on $\mathbf{r}$. Since $\mathbf{r}$ and $t$ are independent, both sides must equal the same constant $E$:

\begin{equation}
i\hbar\,\frac{1}{\phi}\,\frac{d\phi}{dt} = E, \qquad \frac{1}{\psi}\left[-\frac{\hbar^2}{2m}\nabla^2\psi + V\,\psi\right] = E.
\end{equation}

\textbf{Step 4: The Time Equation.}

The first equation gives the time evolution:

\begin{equation}
\frac{d\phi}{dt} = -\frac{iE}{\hbar}\,\phi \quad \Longrightarrow \quad \phi(t) = e^{-iEt/\hbar}.
\end{equation}

This is a pure phase factor with frequency $\omega = E/\hbar$: stationary states oscillate in time with a definite energy $E$. The probability density $|\Psi|^2 = |\psi|^2\,|\phi|^2 = |\psi|^2$ is time-independent---hence the name ``stationary states.''

\textbf{Step 5: The Time-Independent Schr\"odinger Equation.}

The second equation is the eigenvalue equation:

\begin{equation}
-\frac{\hbar^2}{2m}\nabla^2\psi(\mathbf{r}) + V(\mathbf{r})\,\psi(\mathbf{r}) = E\,\psi(\mathbf{r}).
\end{equation}

In operator notation, this is:

\begin{equation}
\boxed{\hat{H}\,\psi = E\,\psi, \qquad \hat{H} = -\frac{\hbar^2}{2m}\nabla^2 + V(\mathbf{r}).}
\end{equation}

This is an eigenvalue equation: $\hat{H}$ is the Hamiltonian operator, $\psi$ is an eigenfunction, and $E$ is the corresponding eigenvalue (energy).

\textbf{Step 6: Physical Content and Boundary Conditions.}

The time-independent Schr\"odinger equation is a second-order linear ODE (in 1D) or PDE (in 3D). Not all values of $E$ yield physically acceptable solutions. The requirements that $\psi$ be:

\begin{itemize}
\item single-valued (uniqueness),
\item continuous (no infinite derivatives),
\item square-integrable ($\int |\psi|^2\,d^3r = 1$, normalizability),
\end{itemize}

restrict $E$ to a discrete spectrum $\{E_n\}$ for bound states (e.g., electrons in atoms) or a continuous spectrum for scattering states. For the hydrogen atom, the discrete energies are $E_n = -13.6\,\text{eV}/n^2$, reproducing the observed spectral lines of hydrogen exactly.

\section*{Characteristics of the Equation}

The equation is a second-order ODE (in 1D) or PDE (in 3D). Solutions exist only for discrete eigenvalues $E_n$ (bound states) or continuous values (scattering states). The eigenfunctions form a complete orthonormal basis. Energy levels scale as $1/n^2$ for hydrogen.
\section*{Solution Methods}

Separation of variables (for separable potentials), power series expansion (hydrogen, harmonic oscillator), numerical methods (finite differences, variational method, perturbation theory, WKB approximation), supersymmetric quantum mechanics.
\section*{Verifications}

Spectral measurements of hydrogen match predictions to $10^{-12}$ precision. Scanning tunneling microscopy images electron wave functions on surfaces. Quantum dots with engineered energy levels produce specific emission wavelengths.
\section*{Applications}

Atomic and molecular physics (energy levels, spectral lines), solid-state physics (band structure, semiconductor device design), quantum chemistry (molecular orbitals), nuclear physics (shell model), quantum computing (qubit energy levels).
\section*{What Richard Feynman Would Have Asked}

\subsection*{1. The Plain-English Translation}
\textit{``What is this equation really asking?''}

The Core Meaning: The Schr\"odinger equation asks: ``For this particular potential energy landscape, what are the allowed wave patterns?'' Just as a guitar string can only vibrate at certain frequencies (harmonics), a quantum particle can only exist in certain wave patterns (eigenstates). Each allowed wave pattern has a specific energy.

The Metaphor: Imagine blowing across the top of bottles of different sizes. Each bottle produces a specific note---you can't get an arbitrary pitch. The Schr\"odinger equation is the same: the ``bottle'' is the potential energy well, and the ``notes'' are the allowed energy levels. Change the bottle shape (the potential), and you change the notes.

\subsection*{2. The Localized Mechanics}
\textit{``What determines the shape of the wave function at a single point?''}

He would press you on the kinetic and potential energy competition. At each point, the kinetic energy term ($-\hbar^2\nabla^2\psi/2m$) tries to smooth out the wave function (reduce curvature), while the potential energy term ($V\psi$) pins the wave function to certain regions. The eigenstate is the balance: a standing wave that satisfies both simultaneously.

The Smoothness Check: He would ask why the wave function must be continuous. Discontinuities would imply infinite kinetic energy (infinite curvature). The only exception is at an infinite potential wall, where the wave function can drop to zero abruptly.

\subsection*{3. Testing the Extreme Limits}
\textit{``What happens when we push the system to extreme conditions?''}

The Infinite Well: For a particle in a box with infinitely high walls, the energy levels are $E_n = n^2 \pi^2 \hbar^2 / (2mL^2)$---they scale as $n^2$. The ground state has nonzero energy ($E_1 > 0$), demonstrating the uncertainty principle: confining the particle ($\Delta x$ small) forces its momentum to be large.

The Hydrogen Atom: The $1/r$ potential gives energy levels $E_n = -13.6 \text{ eV}/n^2$. The wave functions are the famous orbitals ($1s$, $2p$, $3d$, \ldots) whose shapes determine chemical bonding and the periodic table.

The Harmonic Oscillator: The parabolic potential $V = \frac{1}{2}m\omega^2 x^2$ gives equally spaced energy levels $E_n = (n + \frac{1}{2})\hbar\omega$. Even in the ground state, the oscillator has energy $\frac{1}{2}\hbar\omega$---zero-point energy, a purely quantum phenomenon.

\subsection*{4. Dimensionality and Geometry}
\textit{``How does the dimensionality of space change the solutions?''}

He would point out that in 1D, the hydrogen atom doesn't exist (the $1/r$ potential is singular at the origin in 1D). In 2D, hydrogen-like atoms have different energy scaling. In 3D, the familiar $1s$, $2p$, $3d$ orbitals emerge. The angular momentum quantum numbers arise from the angular part of the Laplacian in spherical coordinates.

\subsection*{5. Cross-Disciplinary Connection}
\textit{``Where else does this eigenvalue equation structure appear?''}

The same eigenvalue structure appears in: vibrating drums and strings (normal modes), electromagnetic cavities (resonant frequencies), population ecology (growth modes of interacting species), Google's PageRank algorithm (dominant eigenvector of the web graph), and machine learning (principal component analysis finds the dominant eigenmodes of data). Any system with a linear operator acting on a state vector produces eigenvalue equations.
\section*{FAQ}

\textbf{Q: Why are only discrete energies allowed?}\\
\textbf{A:} Boundary conditions (the wave function must be single-valued, continuous, and normalizable) restrict solutions to specific wavelengths, and since $E$ depends on wavelength, only certain $E$ values are allowed. This is analogous to standing waves on a string: only certain frequencies fit.

\textbf{Q: What happens to the solutions for continuous (unbound) states?}\\
\textbf{A:} For $E > V$ everywhere, the solutions are oscillatory and not normalizable---they represent scattering states (free particles). The energy spectrum is continuous, and the particle is not confined.
\section*{Important Bibliography}

\begin{enumerate}
\item Griffiths, D.J., \textit{Introduction to Quantum Mechanics}, 3rd ed. (Cambridge University Press, 2018), Chapters 2--4.
\item Schr\"odinger, E., ``Quantisierung als Eigenwertproblem,'' \textit{Annalen der Physik} \textbf{79}, 361--376 (1926).
\item Shankar, R., \textit{Principles of Quantum Mechanics}, 2nd ed. (Springer, 2011), Chapters 5--7.
\item Merzbacher, E., \textit{Quantum Mechanics}, 3rd ed. (Wiley, 1998), Chapters 6--8.
\end{enumerate}

